% Generated draft: requires mathematical and source review.
\documentclass[11pt]{article}
\usepackage[margin=25mm]{geometry}
\usepackage{fontspec}
\setmainfont{Latin Modern Roman}
\usepackage{amsmath,amssymb,mathtools,unicode-math}
\setmathfont{Latin Modern Math}
\usepackage{enumitem,needspace,xurl,hyperref}
\setlength{\parindent}{0pt}
\setlength{\parskip}{5pt}
\setlength{\emergencystretch}{3em}
\newcommand{\N}{\mathbb{N}}
\newcommand{\Z}{\mathbb{Z}}
\newcommand{\Q}{\mathbb{Q}}
\newcommand{\R}{\mathbb{R}}
\newcommand{\C}{\mathbb{C}}
\newcommand{\F}{\mathbb{F}}
\newcommand{\A}{\mathbb{A}}
\newcommand{\PP}{\mathbb{P}}
\newcommand{\E}{\mathbb{E}}
\newcommand{\eps}{\varepsilon}
\newcommand{\dd}{\,\mathrm{d}}
\begin{document}
\section*{UnsolvedMath ID: 30001612}
\textbf{Source problem code:} catalogue entry
\subsection{Definitions and mathematical statement}
Let \(\Omega\subset\C\) be a domain described in the supplied source as a ``crescent domain'' or ``crescent-type domain,'' with bounded complementary component \(G\) and unbounded complementary component \(\Omega_\infty\). The source record does not define the crescent condition or specify the boundary regularity it entails; that geometric qualification must therefore be retained pending clarification. Set
\[
\delta(z)=\operatorname{dist}(z,\Omega_\infty)
=\inf_{w\in\Omega_\infty}|z-w|.
\]
Let \(\omega\) be the source's harmonic measure on \(\partial G\). Its defining domain, side of the boundary, and pole are not specified in the supplied record.

Write \(dA\) for planar area measure and, for \(1\leq p<\infty\), define
\[
L_a^p(\Omega,dA)=\left\{f:\Omega\to\C:\ f\text{ is holomorphic and }\int_\Omega |f(z)|^p\,dA(z)<\infty\right\}.
\]
The source asks whether
\[
\int_{\partial G}\log\delta(z)\,d\omega(z)=-\infty
\quad\Longrightarrow\quad
H^p(\Omega,dA)=L_a^p(\Omega,dA)
\quad\text{for every }1\leq p<\infty.
\]
Here \(\log 0=-\infty\), and the integral condition is understood as an extended Lebesgue integral equal to \(-\infty\), rather than an undefined difference of infinite quantities.

The supplied record does not define \(H^p(\Omega,dA)\). Its title and literature assessment indicate that polynomial density is intended. To express that interpretation explicitly, put
\[
\mathcal P_p(\Omega)=\{P|_\Omega:P\in\C[z],\ P|_\Omega\in L^p(\Omega,dA)\}.
\]
If the intended definition is
\[
H^p(\Omega,dA)=\overline{\mathcal P_p(\Omega)}^{\,L^p(\Omega,dA)},
\]
then the proposed conclusion says that, for every real \(p\) with \(1\leq p<\infty\), every \(f\in L_a^p(\Omega,dA)\), and every \(\varepsilon>0\), there exists a polynomial \(P\) whose restriction belongs to \(L^p(\Omega,dA)\) such that
\[
\left(\int_\Omega |f(z)-P(z)|^p\,dA(z)\right)^{1/p}<\varepsilon.
\]
This interpretation of \(H^p\) is conditional and requires confirmation from the source, together with the precise crescent-domain and harmonic-measure conventions.
\subsection{Short English statement}
The question asks whether divergence to \(-\infty\) of the harmonic-measure integral of the logarithmic distance to the outer complementary component forces polynomial density in every finite-exponent Bergman space. The crescent geometry, harmonic measure, and notation \(H^p\) require clarification before this interpretation is fully specified.
\subsection{Sources}
\begin{itemize}
\item Local UnsolvedMath registry: \texttt{data/unsolvedmath/problems.json}, permanent ID \(30001612\), source problem code \texttt{catalogue entry}, title ``Polynomial Density on Crescent Domains,'' in the Oberwolfach Reports Open Problems collection. No individual proposer is identified in the supplied record.
\item Supplied report citation: \emph{Operator Theory and Harmonic Analysis} (2011), Oberwolfach Reports, pp.~2813--2875, DOI \texttt{10.4171/owr/2010/49}. \url{https://doi.org/10.4171/owr/2010/49}. The record assigns proposed year 2010. Its \texttt{corrected\_verified} label and assertion that the statement was checked against the official report are supplied verification metadata, not an independent verification here.
\item Supplied literature reference: ``Hedenmalm and Shimorin weighted Bergman estimates,'' \url{https://arxiv.org/abs/math/0406345}. No fuller bibliographic details or theorem locator are supplied.
\item Supplied literature reference: ``original problem report,'' \url{https://doi.org/10.4171/owr/2010/49}.
\item Literature checked at: \texttt{2026-08-21}. This is supplied dated curation, not independently verified. The supplied status is ``partially solved.'' The accompanying assessment states that the divergence condition implies polynomial density for \(1\leq p<3+\tau\), for a universal constant \(\tau>0\), while the assertion for every finite \(p\) remains unresolved. The background describes this as dated literature triage derived from the supplied Oberwolfach curation and requests independent verification of time-sensitive status and cited claims.
\end{itemize}
\end{document}
